\section{Implicaciones prácticas: cómo afecta al entrenamiento la diferencia de distribución}
Una vez que se entiende SFT y RL desde la perspectiva de la distribución, resulta más fácil ver por qué ambos entrenan modelos con estilos distintos. SFT se parece más a aproximar la distribución empírica de los datos existentes, mientras que RL empuja activamente la masa de probabilidad hacia regiones de alta recompensa; por eso los compromisos entre diversidad, estabilidad y controlabilidad también son distintos entre ambos.

\begin{knowledgebox}{Qué preguntar primero al experimentar}
\begin{itemize}
  \item La tarea actual, ¿necesita más apegarse a la distribución humana o más llevar la recompensa al extremo
  \item ¿El rollout es suficientemente en línea como para sostener el desplazamiento de distribución de RL
  \item ¿Es necesario primero sentar una base con SFT y después usar RL para un refuerzo localizado
\end{itemize}
\end{knowledgebox}

\subsection{Resumen de la sección}
El mayor valor de la perspectiva de distribución está en que ayuda a quien investiga a juzgar cuándo imitar y cuándo buscar la recompensa, y cómo deberían combinarse ambos enfoques.

\newpage
